% !TeX program = lualatex
% halo:begin
% {
%   "version": 1,
%   "publish": true,
%   "course": "mathematical-analysis-i",
%   "section": "3.2",
%   "title": "数学分析（一）3.2：函数极限的性质",
%   "categories": ["study-notes"],
%   "tags": ["mathematics"],
%   "excerpt": "函数极限的唯一性、局部有界性、保号性、保序性、夹逼定理、四则运算法则及归结原则。"
% }
% halo:end
% 来源：IMG_2167.HEIC—IMG_2170.HEIC，并与 IMG_2165.HEIC—IMG_2166.HEIC 中的复习内容对照。
% 归结原则的逆向构造补全了板书未写完的证明；记号沿用 03-01。
\RequirePackage{iftex}
\RequireLuaTeX
\documentclass[UTF8,fontset=fandol,a4paper,12pt]{ctexart}
\usepackage[margin=2.5cm]{geometry}
\usepackage{amsmath,amssymb,amsthm}
\usepackage{enumitem,fancyhdr}
\usepackage[hidelinks]{hyperref}
\newcommand{\R}{\mathbb R}
\newcommand{\Nplus}{\mathbb N_{+}}
\newcommand{\U}{U}
\theoremstyle{definition}
\newtheorem*{definition}{定义}
\newtheorem{theorem}{定理}
\renewcommand{\thetheorem}{3.\arabic{theorem}}
\newtheorem*{example}{例题}
\renewcommand{\proofname}{证明}
\setlist[enumerate]{leftmargin=2em,itemsep=0.2em,topsep=0.3em}
\setlength{\headheight}{16pt}
\pagestyle{fancy}
\fancyhf{}
\fancyhead[L]{\small 数学分析（一）}
\fancyhead[R]{\small 3.2\quad 函数极限的性质}
\fancyfoot[C]{\thepage}
\renewcommand{\headrulewidth}{0.3pt}
\raggedbottom
\title{函数极限的性质}
\author{Yuchen Zhang}
\date{整理日期：2026年10月14日}
\makeatletter
\renewcommand{\maketitle}{\begin{center}
{\LARGE\bfseries \@title\par}\vspace{0.6em}
\@author\par\vspace{0.2em}{\small\@date\par}
\end{center}}
\makeatother
\begin{document}
\maketitle
\setcounter{section}{2}
\setcounter{theorem}{0}
\section{函数极限的性质}

以下各定理均在同一种趋近方式下讨论，例如 $x\to x_0$、$x\to x_0^+$ 或 $x\to+\infty$。为简洁起见，正文先写 $x\to x_0$；单侧极限和无穷远处的结论完全相同，只需把邻域条件换成相应的单侧或远端条件。

\subsection*{唯一性}
\begin{theorem}[唯一性]
若
\[
 \lim_{x\to x_0}f(x)=A,\qquad \lim_{x\to x_0}f(x)=B,
\]
则 $A=B$。
\end{theorem}
\begin{proof}
反设 $A\ne B$，取 $\varepsilon=|A-B|/3>0$。由两个极限的定义，存在 $\delta_1,\delta_2>0$，使得当 $0<|x-x_0|<\delta_i$ 时，分别有
\[
 |f(x)-A|<\varepsilon,\qquad |f(x)-B|<\varepsilon.
\]
取 $\delta=\min\{\delta_1,\delta_2\}$，任取 $0<|x-x_0|<\delta$，则
\[
 |A-B|\le |A-f(x)|+|f(x)-B|<2\varepsilon
 =\frac23|A-B|,
\]
矛盾。故 $A=B$。
\end{proof}

\subsection*{局部有界性}
\begin{theorem}[局部有界性]
若 $\lim_{x\to x_0}f(x)=A$，则存在 $\delta>0$ 和 $M>0$，使得当
\[
 0<|x-x_0|<\delta
\]
时有 $|f(x)|<M$。
\end{theorem}
\begin{proof}
在极限定义中取 $\varepsilon=1$，可得某个 $\delta>0$，使
\[
 |f(x)-A|<1
\]
对所有 $0<|x-x_0|<\delta$ 成立。因此
\[
 |f(x)|\le |f(x)-A|+|A|<1+|A|.
\]
取 $M=1+|A|$ 即可。
\end{proof}

这里的“局部”是指某个去心邻域；极限存在并不要求函数在整个定义域上有界。局部有界性本身也不能推出极限存在，例如 $\sin(1/x)$ 在 $x\to0$ 时有界但没有极限。

\subsection*{保号性}
\begin{theorem}[保号性]
若 $\lim_{x\to x_0}f(x)=A>0$，则存在 $\delta>0$，使 $0<|x-x_0|<\delta$ 时 $f(x)>0$。若 $A<0$，则在某个去心邻域内 $f(x)<0$。
\end{theorem}
\begin{proof}
设 $A>0$。取 $\varepsilon=A/2$，则存在 $\delta>0$，使
\[
 |f(x)-A|<\frac A2.
\]
于是
\[
 f(x)>A-\frac A2=\frac A2>0
\]
对所有 $0<|x-x_0|<\delta$ 成立。

当 $A<0$ 时，取 $\varepsilon=|A|/2$，可得
\[
 f(x)<A+\frac{|A|}{2}=\frac A2<0.
\]
\end{proof}

保号性只适用于极限严格大于零或严格小于零的情形。若极限为零，函数值可以在邻域内变号，例如 $f(x)=x$ 在 $x\to0$ 时既有正值也有负值。

\subsection*{保序性与夹逼定理}
\begin{theorem}[保序性]
若 $\lim_{x\to x_0}f(x)=A$、$\lim_{x\to x_0}g(x)=B$，并且存在 $\delta_0>0$，使
\[
 f(x)\le g(x)\qquad(0<|x-x_0|<\delta_0),
\]
则 $A\le B$。
\end{theorem}
\begin{proof}
反设 $A>B$。取 $\varepsilon=(A-B)/3$。当 $x$ 充分接近 $x_0$ 时，
\[
 f(x)>A-\varepsilon>B+\varepsilon>g(x),
\]
与 $f(x)\le g(x)$ 矛盾。因此 $A\le B$。
\end{proof}

\begin{theorem}[夹逼定理]
若在某个去心邻域内有
\[
 f(x)\le h(x)\le g(x),
\]
且
\[
 \lim_{x\to x_0}f(x)=\lim_{x\to x_0}g(x)=A,
\]
则 $\lim_{x\to x_0}h(x)=A$。
\end{theorem}
\begin{proof}
任取 $\varepsilon>0$。取 $\delta>0$，使 $0<|x-x_0|<\delta$ 时既有夹逼关系，又有
\[
 |f(x)-A|<\varepsilon,\qquad |g(x)-A|<\varepsilon.
\]
于是
\[
 A-\varepsilon<f(x)\le h(x)\le g(x)<A+\varepsilon,
\]
即 $|h(x)-A|<\varepsilon$，故 $h(x)\to A$。
\end{proof}

\begin{example}
证明
\[
 \lim_{x\to0}x\left\lfloor\frac1x\right\rfloor=1.
\]
\end{example}
\begin{proof}
先设 $x>0$。由取整函数的性质，
\[
 \frac1x-1<\left\lfloor\frac1x\right\rfloor\le\frac1x.
\]
乘以正数 $x$，得
\[
 1-x<x\left\lfloor\frac1x\right\rfloor\le1.
\]
当 $x\to0^+$ 时，由夹逼定理，右侧函数趋于 $1$。当 $x<0$ 时，同样由
\[
 \frac1x-1<\left\lfloor\frac1x\right\rfloor\le\frac1x
\]
乘以负数 $x$ 得
\[
 1\le x\left\lfloor\frac1x\right\rfloor<1-x,
\]
故右侧也趋于 $1$。再由双侧极限判定，原极限为 $1$。
\end{proof}

\subsection*{四则运算法则}
\begin{theorem}[四则运算法则]
若 $\lim_{x\to x_0}f(x)=A$、$\lim_{x\to x_0}g(x)=B$，则
\[
 \lim_{x\to x_0}[f(x)\pm g(x)]=A\pm B,
\]
\[
 \lim_{x\to x_0}f(x)g(x)=AB.
\]
若 $B\ne0$，则在某个去心邻域内 $g(x)\ne0$，且
\[
 \lim_{x\to x_0}\frac{f(x)}{g(x)}=\frac AB.
\]
\end{theorem}
\begin{proof}
加减法由
\[
 |[f(x)\pm g(x)]-(A\pm B)|
 \le |f(x)-A|+|g(x)-B|
\]
及 $\varepsilon/2$ 的选取立即得到。

乘法中，先由局部有界性取 $M>0$，使 $|f(x)|<M$。于是
\[
 |f(x)g(x)-AB|
 \le |f(x)|\,|g(x)-B|+|B|\,|f(x)-A|
 <M|g(x)-B|+|B|\,|f(x)-A|.
\]
右端两项可同时取得任意小，故乘积极限为 $AB$。

若 $B\ne0$，由 $g(x)\to B$，取 $\varepsilon=|B|/2$，可得
\[
 |g(x)|>\frac{|B|}{2}>0
\]
在某个去心邻域内成立。因此
\[
 \left|\frac1{g(x)}-\frac1B\right|
 =\frac{|g(x)-B|}{|B|\,|g(x)|}
 \le\frac2{|B|^2}|g(x)-B|\to0.
\]
再用乘法法则即得商的结论。
\end{proof}

\begin{example}
计算
\[
 \lim_{x\to+\infty}\frac{2x^2+3x}{x^2+2x+1}.
\]
\end{example}
\begin{proof}[解]
除以 $x^2$ 后为
\[
 \frac{2+3/x}{1+2/x+1/x^2}\longrightarrow2.
\]
\end{proof}

\subsection*{归结原则}
直接从函数极限定义判断，有时需要同时处理所有充分接近 $x_0$ 的点。归结原则把它转化为数列问题。

\begin{theorem}[归结原则（Heine 定理）]
设 $f$ 在 $x_0$ 的某个去心邻域 $\U^0(x_0,\delta_0)$ 内有定义，则
\[
 \lim_{x\to x_0}f(x)=A
\]
当且仅当对任意满足
\[
 x_n\in\U^0(x_0,\delta_0),\qquad x_n\to x_0
\]
的数列 $\{x_n\}$，都有
\[
 f(x_n)\to A.
\]
\end{theorem}
\begin{proof}
先证必要性。设 $f(x)\to A$，任取 $x_n\in\U^0(x_0,\delta_0)$ 且 $x_n\to x_0$。给定 $\varepsilon>0$，由函数极限定义取 $\delta>0$，使
\[
 0<|x-x_0|<\delta\Longrightarrow |f(x)-A|<\varepsilon.
\]
又由 $x_n\to x_0$，存在 $N$，使 $n>N$ 时 $|x_n-x_0|<\delta$。因 $x_n\ne x_0$，故 $|f(x_n)-A|<\varepsilon$，即 $f(x_n)\to A$。

再证充分性。反设函数极限不等于 $A$。则存在 $\varepsilon_0>0$，使任意 $\delta>0$ 都能找到满足 $0<|x-x_0|<\delta$ 的点 $x$，使
\[
 |f(x)-A|\ge\varepsilon_0.
\]
对每个 $n\in\Nplus$ 取 $x_n$ 满足
\[
 0<|x_n-x_0|<\min\left\{\delta_0,\frac1n\right\},\qquad |f(x_n)-A|\ge\varepsilon_0.
\]
则 $x_n\in\U^0(x_0,\delta_0)$ 且 $x_n\to x_0$，但 $f(x_n)$ 不趋于 $A$，与假设矛盾。因此 $\lim_{x\to x_0}f(x)=A$。
\end{proof}

\begin{example}
证明 $\displaystyle\lim_{x\to0}\sin(1/x)$ 不存在。
\end{example}
\begin{proof}
取
\[
 x_n=\frac1{2\pi n+\pi/2},\qquad
 y_n=\frac1{2\pi n+3\pi/2}.
\]
两列都趋于 $0$，但
\[
 \sin\frac1{x_n}=1,\qquad \sin\frac1{y_n}=-1.
\]
若 $\lim_{x\to0}\sin(1/x)$ 存在，则由归结原则，两列函数值应趋于同一个极限，矛盾。故该极限不存在。
\end{proof}

\begin{example}
计算 $\displaystyle\lim_{x\to\pi/4}(x\tan x-1)$。
\end{example}
\begin{proof}[解]
由函数极限的四则运算法则和三角函数的连续性，
\[
 \lim_{x\to\pi/4}x\tan x
 =\frac\pi4\cdot\tan\frac\pi4=\frac\pi4.
\]
因此
\[
 \lim_{x\to\pi/4}(x\tan x-1)=\frac\pi4-1.
\]
\end{proof}

\end{document}
